\clearpage
\phantomsection% 
\addcontentsline{toc}{chapter}{KATA PENGANTAR}
%\thispagestyle{fancy}

\begin{justifying}
	\large \bfseries \centering \MakeUppercase{Kata Pengantar}\linebreak
	
	\normalsize \normalfont \justifying
	Puji syukur kehadirat Tuhan Yang Maha Esa, yang atas berkat rahmat dan karunia-Nya, penulis dapat menyelesaikan tugas akhir dan menyusun laporan dengan baik dan tepat waktu. Tugas akhir dengan judul ``{\thetitle}'' disusun sebagai salah satu syarat untuk memenuhi persyaratan kelulusan program studi S-1 Teknik Informatika di Institut Teknologi Sumatera. \par
	
	Penulis menyadari bahwa kelancaran pelaksanaan penelitian serta penyusunan tugas akhir ini tidak terlepas dari bimbingan, dukungan, doa, dan semangat dari berbagai pihak. Oleh karena itu, dengan segala kerendahan hati, penulis ingin mengucapkan terima kasih yang sebesar-besarnya kepada: \par
	\begin{enumerate}
		\item Kedua orang tua, Ayah dan Ibu. Terima kasih atas segala curahan kasih sayang, dukungan moril dan materiil, serta doa yang tiada pernah putus yang senantiasa mengiringi setiap langkah penulis.
		\item Bapak Meida Cahyo Untoro, S.Kom., M.Kom. selaku Dosen Pembimbing Tugas Akhir yang telah dengan sabar memberikan waktu, masukan, arahan, dan bimbingan dalam proses penelitian dan penyusunan tugas akhir ini.
		\item Teman-teman dan rekan-rekan penulis selama masa perkuliahan. Terima kasih atas semangat, kerja sama, dan diskusi yang membangun selama ini.
	\end{enumerate} \par
	Penulis menyadari sepenuhnya bahwa tugas akhir ini masih memiliki banyak kekurangan karena keterbatasan pengalaman dan pengetahuan. Oleh karena itu, penulis sangat mengharapkan kritik dan saran yang membangun dari semua pihak demi perbaikan di masa mendatang. \par
	
	Akhir kata, semoga tugas akhir ini dapat memberikan manfaat, baik bagi penulis pribadi maupun bagi pembaca.
	\vfill
	
\end{justifying}
\clearpage





\begin{comment}
\clearpage
\pagestyle{fancy}
\fancyhf{}
\fancyhead[R]{\thepage}
\phantomsection% 
%\clearpage
%\phantomsection% 
\addcontentsline{toc}{chapter}{KATA PENGANTAR}
%\thispagestyle{fancy}

\begin{justifying}
	\large \bfseries \centering \MakeUppercase{Kata Pengantar}\linebreak
	
	\normalsize \normalfont \justifying
	Puji syukur kehadirat Tuhan Yang Maha Esa atas limpahan rahmat, karunia, serta petunjuk-Nya sehingga penyusunan tugas akhir ini telah terselesaikan dengan baik. Dalam penyusunan tugas akhir ini penulis telah banyak mendapatkan arahan, bantuan, serta dukungan dari berbagai pihak. Oleh karena itu pada kesempatan ini penulis mengucapan terima kasih kepada: \par
	\begin{enumerate}
		\item Bapak Prof. Dr. I. Nyoman Pugeg Aryantha selaku Rektor Institut Teknologi Sumatera.  
		\item Bapak Hadi Teguh Yudistira, S.T., Ph.D. selaku Dekan Fakultas Teknologi Industri.
		\item Bapak Andika Setiawan, S. Kom., M. Cs. selaku Ketua Program Studi Teknik Informatika.
		\item Bapak Ilham Firman Ashari, S. Kom., M.T. selaku Koordinator Tugas Akhir Program Studi Teknik Informatika.  
		\item Bapak Martin C. T. Manullang, Ph.D. selaku Dosen Pembimbing atas ide, waktu, tenaga, perhatian, dan masukan yang telah disumbangsihkan kepada penulis.
            \item Bapak Andika Setiawan, S. Kom., M. Cs dan Bapak Eko Dwi Nugroho, S.Kom., M.Cs. selaku Dosen Penguji atas saran dan masukan yang diberikan. 
		\item Kedua Orang Tua dan Adik yang selalu memberikan dukungan dan doa sehingga penulis dapat menyelesaikan tugas akhir ini. Kelembutan, dukungan, dan cinta yang kalian berikan selalu menjadi sumber inspirasi dan kekuatan.
		\item Teman-teman penulis yang membantu selama masa perkuliahan yang tidak bisa disebutkan satu persatu.
	\end{enumerate} \par
	Akhir kata penulis berharap semoga tugas akhir ini dapat memberikan manfaat bagi kita semua. Penulis menyadari bahwa tugas akhir ini tidak luput dari kekurangan dan kelemahan, dan penulis  terbuka untuk menerima saran, kritik, dan masukan yang membangun.
	\vfill
	
\end{justifying}
\clearpage
\end{comment}